\documentclass{article}

% Packages
\usepackage{amsmath, amsfonts, amssymb, amsthm, bm, bbm} 
\usepackage{enumerate, graphicx, xcolor, tcolorbox, hyphenat, circledsteps}
\usepackage[shortlabels]{enumitem}
\usepackage{url, hyperref}
\hypersetup{
	colorlinks   = true,
	urlcolor     = cyan,
	linkcolor    = red,
	citecolor    = black
}

% Formatting
\usepackage[a4paper, noheadfoot, scale={0.7,0.8}]{geometry}
\title{Wk 2 Supervision Problems 1}
\author{Luca Bollini}
\date{13 October 2026}
\hyphenpenalty=10000

% Commands
\newcommand{\N}{\mathbb{N}}
\newcommand{\Z}{\mathbb{Z}}
\newcommand{\Q}{\mathbb{Q}}
\newcommand{\R}{\mathbb{R}}
\newcommand{\ddt}[1]{\frac{d}{d#1}}
\newcommand{\ddttwo}[1]{\frac{d^2}{d#1^2}}
\newcommand{\ts}{\thinspace}

% Colours
\newtcolorbox{questionbox}{
	colframe=black, 
	colback=white, 
	boxrule=0.2mm,
	arc=0mm,
}


\begin{document}
\maketitle

% LESSON PLAN: Split the group in two and have half of them work on problem A on the left side of the blackboard, while the others work on problem B on the right side of the blackboard.

% NOTE: This is meant to be an exercise in deciphering word problems and turning them into maths, similar to the first modelling assignment.

\begin{center}
\huge{\textbf{A}} 
\end{center}

Newton's law of cooling states that the rate an object cools is proportional to the temperature difference between the object and its surroundings.
A $90^\circ$C cup of coffee is placed in a room maintained at $20^\circ$C. After 10 minutes, the coffee has cooled to $70^\circ$C.

\begin{enumerate}
    \item Set up and solve an initial value problem for the temperature, $T(t)$, of the coffee using Newton's law of cooling.
    \item Find the time at which the coffee reaches $40^\circ$C.
\end{enumerate}

\vspace{5pt}

\begin{center}
\huge{\textbf{B}}
\end{center}

The displacement of a mass on a spring can be described by Newton's second law as follows:
\[
m\ddttwo{t}x(t) + c\ddt{t}x(t) + kx(t) = 0
\]
where $m$ is the mass, $c$ is the damping constant, and $k$ is the spring constant.
A mass of $1$ kg is attached to a spring with constant $k=13$ N/m. The system experiences damping with coefficient $c=4$ Ns/m. At $t=0$ the mass is displaced $0.5$ m from equilibrium and released from rest.

\begin{enumerate}
    \item Write down the initial value problem governing the displacement of the mass, $x(t)$, and find the complimentary function.
    \item Solve the problem completely and determine the long-time behaviour of the system.
\end{enumerate}
\end{document}
